\section*{अध्याय \ref{ch:b3:genetic-engineering} --- आनुवंशिक अभियांत्रिकी और जैवप्रौद्योगिकी}

\begin{solution}{exo:b3:genetic-engineering:1}
प्रतिकृति का कोई उद्गम, ताकि परपोषी \omterm{def:b3:genetic-engineering:tools}{प्लाज़्मिड} की नकल करे और उसे पुत्री
कोशिकाओं तक पहुँचाए; कोई \omterm{def:b3:genetic-engineering:tools}{वरणीय चिह्नक}, प्रायः प्रतिजैविक-प्रतिरोध,
ताकि केवल वही कोशिकाएँ उगें जो \omterm{def:b3:genetic-engineering:tools}{प्लाज़्मिड} ढो रही हैं; और कोई अनन्य
प्रतिबंधन स्थल (या स्थलों का गुच्छा) जिस पर अंतर्वेशन जोड़ा जाए, बिना
\omterm{def:b3:genetic-engineering:tools}{प्लाज़्मिड} को कहीं और काटे।
\end{solution}

\begin{solution}{exo:b3:genetic-engineering:2}
यादृच्छिक DNA में कोई दिया हुआ छह-क्षार अनुक्रम हर $4^{6} = 4096$ bp में
एक बार आता है: लगभग $4.6\times 10^{6}/4096 \approx 1100$ खंड, औसतन
\qty{4.1}{kb} के। असली जीनोम यादृच्छिक नहीं होते --- क्षार-संघटन,
कूटक-उपयोग और कुछ अनुक्रमों से बचाव गिनती खिसका देते हैं
(\emph{E.~coli} में EcoRI स्थलों की वास्तविक संख्या कुछ सौ है)।
\end{solution}

\begin{solution}{exo:b3:genetic-engineering:3}
\qty{95}{\celsius} पर विकृतीकरण लड़ियाँ अलग करता है;
\qtyrange{50}{65}{\celsius} पर तापानुशीतन \omterm{def:b3:genetic-engineering:pcr}{प्रारंभकों} को उनके स्थलों से
युग्मित होने देता है; \qty{72}{\celsius} पर दीर्घन पॉलीमरेज़ को हर
\omterm{def:b3:genetic-engineering:pcr}{प्रारंभक} से नकल करने देता है। कोई साधारण पॉलीमरेज़ पहले ही चक्र में
\qty{95}{\celsius} पर नष्ट हो जाती और उसे तीस बार नए सिरे से जोड़ना
पड़ता; Taq सारे चक्र झेल जाती है।
\end{solution}

\begin{solution}{exo:b3:genetic-engineering:4}
कोई ऐसा \omterm{def:b3:genetic-engineering:crispr}{मार्गदर्शक RNA} जिसके बीस क्षार लक्ष्य से मेल खाएँ, और
अलक्ष्य लड़ी पर लक्ष्य के ठीक 3$'$ ओर कोई \omterm{def:b3:genetic-engineering:crispr}{PAM} (NGG)। \omterm{def:b3:dna-repair:lesions}{द्विलड़ी विखंडन} के
बाद: सिरा-संयोजन, जो कोई छोटा अंतर्वेशन या विलोपन छोड़ता है और प्रायः
जीन का वाचन-ढाँचा नष्ट कर देता है; या वांछित अनुक्रम ढोते किसी साँचे
के साथ \omterm{def:b3:dna-repair:dsb}{समजात पुनर्संयोजन}, जो उसे लिख देता है।
\end{solution}

\begin{solution}{exo:b3:genetic-engineering:5}
$50\times 1.9^{35} = 50\times 5.7\times 10^{9} \approx 3\times 10^{11}$
प्रतियाँ। $\Delta C_{T} = 7.3$; अनुपात $1.95^{7.3} = e^{7.3\ln 1.95}
\approx 130$: पहले प्रतिदर्श में लगभग $130$ गुना अधिक साँचा था।
\end{solution}

\begin{solution}{exo:b3:genetic-engineering:6}
जीवाणु विच्छेदन नहीं कर सकते: अंतर्वेशनों वाली कोई जीनोमी प्रति ऐसे
संदेश में अनुलेखित होती जो काम का न होता, इसलिए परिपक्व mRNA से नकल की
गई, अंतर्वेशन-रहित cDNA काम में ली जाती है। दूसरी अड़चनें: मानव प्रवर्तक
को जीवाणु पॉलीमरेज़ नहीं पहचानती (कोई जीवाणु प्रवर्तक और
राइबोसोम-बंधक स्थल दीजिए); मानव कूटक-उपयोग भिन्न है (जीन को परपोषी के
पसंदीदा कूटकों से संश्लेषित कीजिए); प्रोटीन वलित न हो या पुंजित हो जाए
(ताप घटाइए, किसी विलेय साथी से संलयित कीजिए, या अंतर्वेशी पिंडों से
फिर वलित कीजिए); ग्लाइकोसिलीकरण नहीं होता (यदि वह मायने रखे तो खमीर या
स्तनधारी कोशिकाएँ लीजिए); डाइसल्फ़ाइडों को ऑक्सीकारी परिकोशिकाद्रव्य चाहिए।
\end{solution}

\begin{solution}{exo:b3:genetic-engineering:7}
लक्षित स्तंभ कोशिकाएँ किसी परपोषी कंदुक में अंतःक्षेपित की जाती हैं और
उसकी अपनी कोशिकाओं से मिल जाती हैं, इसलिए चूहा दो जीन-प्ररूपों से बना
होता है --- एक विचित्रकाय --- और वह युग्मविकल्पी तभी आगे देता है जब
उसकी कुछ जनन-कोशिकाएँ लक्षित कोशिकाओं से बनी हों। विचित्रकाय का
वन्य-प्ररूप चूहे से प्रजनन विषमयुग्मजी देता है (यदि जनन-वंश आधा लक्षित
हो तो उसके आधे युग्मकों से); विषमयुग्मजियों का अंतःसंकरण एक-चौथाई
समयुग्मजी \omterm{def:b3:genetic-engineering:transgenic}{निष्क्रियण} देता है।
\end{solution}

\begin{solution}{exo:b3:genetic-engineering:8}
ठीक-ठीक: $6.4\times 10^{9}\times 4^{-20}\times (1/16) \approx 4\times
10^{-4}$ स्थल (\omterm{def:b3:genetic-engineering:crispr}{PAM} के GG की कीमत $16$ का गुणक है) --- व्यवहार में कोई
नहीं। दो असंगतियों के भीतर $1 + 60 + 1710 = 1771$
अनुक्रम हैं: $1771\times 4\times 10^{-4} \approx 0.7$ संयोगी स्थल।
असली जीनोम यादृच्छिक नहीं है, और किसी मार्गदर्शक की जाँच सीधे उसी के
सामने की जाती है।
\end{solution}

\begin{solution}{exo:b3:genetic-engineering:9}
$p_{1} = 0.05 + 0.8\times 0.05\times 0.95 = 0.088$; $p_{2} = 0.152$;
$p_{3} = 0.255$; फिर $0.41$, $0.60$: $0.5$ पार करने में लगभग पाँच
पीढ़ियाँ, आकलन $\ln(1/(2\times 0.05))/\ln 1.8 \approx 4$ के सामने।
\end{solution}

\begin{solution}{exo:b3:genetic-engineering:10}
सिरा-संयोजन हर कोशिका में चलता है, उन शांत स्तंभ कोशिकाओं में भी जिनमें
\omterm{def:b3:dna-repair:dsb}{समजात पुनर्संयोजन}, जो S/G2 की प्रक्रिया है, अदक्ष है; उसे कोई साँचा
पहुँचाना नहीं पड़ता और उसका उत्पाद --- संवर्धक का कोई भी तोड़ --- काम
पूरा कर देता है, जबकि सुधार के लिए ठीक-ठीक अनुक्रम चाहिए और वह
कोशिकाओं का अल्पांश ही देता है। साथ ही, वही संपादन हर $\beta$-ग्लोबिन
उत्परिवर्तन के लिए काम करता है, दात्र हो या थैलेसीमिक। हानि: वह कोई
नियामक अवयव हटा देता है (उसकी और जो भी भूमिकाएँ हों उनके साथ), दात्र
युग्मविकल्पी अपनी जगह छोड़ जाता है, और अंतर्वेशन-विलोपनों का मिश्रण
अनियंत्रित रहता है।
\end{solution}

\begin{solution}{exo:b3:genetic-engineering:11}
समयुग्मजी अनुकूलता $1 - s$ के साथ: $p' = \bigl[p^{2}(1-s) + p(1-p)(1+c)
\bigr]/(1 - sp^{2})$। दुर्लभता पर समयुग्मजी नगण्य हैं और चालक
$(1+c)p$ की तरह बढ़ता है, $s$ जो भी हो; वह स्थिरीकरण तक जाएगा या नहीं, यह $p = 1$
के पास तय होता है, जहाँ बारंबारता $q$ वाला कोई दुर्लभ वन्य-प्ररूप
युग्मविकल्पी $q' \approx q(1-c)/(1-s)$ के रूप में लौटता है: $c > s$ हो
तो स्थिरीकरण, अन्यथा कोई मध्यवर्ती साम्य। प्रतिरोध: हर विफल रूपांतरण
(पुनर्संयोजन के बजाय सिरा-संयोजन) ऐसा युग्मविकल्पी छोड़ जाता है जिसे
मार्गदर्शक अब नहीं पहचानता; यदि ऐसे युग्मविकल्पी अनुकूल हों --- किसी
अनावश्यक क्षेत्र में कोई ढाँचा-रक्षी अंतर्वेशन-विलोपन --- तो उन पर कोई
लागत नहीं जबकि चालक पर है, इसलिए चयन उन्हें वरीयता देता है और वे चालक
की जगह ले लेते हैं। उपाय हैं ऐसे अनिवार्य अनुक्रमों को लक्ष्य बनाना
जहाँ अंतर्वेशन-विलोपन घातक हों, और एक साथ कई मार्गदर्शक काम में लेना।
\end{solution}

\begin{solution}{exo:b3:genetic-engineering:12}
किसी युग्मनज में अंतःक्षेपित \omterm{def:b3:genetic-engineering:crispr}{Cas9} पहले विभाजनों के बाद भी काम कर सकती
है, इसलिए भिन्न कोशिकाओं को भिन्न मरम्मत मिलती है --- विचित्रता; हर
विखंडन की मरम्मत कोशिका चुनती है, और सिरा-संयोजन कटाव पर अप्रत्याशित
अंतर्वेशन-विलोपन तथा बड़े विलोपन या विषमयुग्मजता-ह्रास देता है;
समजात-पुनर्संयोजन परिणाम अल्पांश होता है; लक्ष्येतर विखंडन ऐसे स्थलों
पर होते हैं जिनका पूर्वानुमान पूरा नहीं लगाया जा सकता; और भ्रूण की
जाँच केवल कुछ बायोप्सी कोशिकाओं के अनुक्रमण से हो सकती है, जो शेष के
लिए नहीं बोल सकतीं। जो प्रमाण चाहिए: ऐसी विधियाँ जो हर कोशिका को एक
जैसा संपादित करें (पहली S प्रावस्था से पहले) और जिनकी दक्षता
\qty{100}{\%} के निकट हो, संपादित पशु-भ्रूणों की हर कोशिका का
पूर्ण-जीनोम अनुक्रमण जो कोई अनपेक्षित बदलाव न दिखाए, और पीढ़ियों तक
संपादित पशुओं का दीर्घकालिक अनुवर्तन --- इनमें से कुछ भी उपलब्ध नहीं।
\end{solution}

\begin{solution}{pb:b3:genetic-engineering:1}
\textbf{1.} प्रति $4096$ bp एक बार। $P(\text{कोई स्थल नहीं}) = (1 - 1/4096)^{1500}
= e^{-0.37} = 0.69$।
\textbf{2.} ऐसी एंज़ाइम चुनिए जिसका जीन में कोई स्थल न हो; या जीन को
PCR से ऐसे \omterm{def:b3:genetic-engineering:pcr}{प्रारंभकों} के साथ प्रवर्धित कीजिए जो अपने 5$'$
सिरों पर नए प्रतिबंधन स्थल ढोते हों (या भोथरा अथवा जोड़-निरपेक्ष
क्लोनन काम में लीजिए)।
\textbf{3.} \qty{0.1}{\micro g}\,$\times 2\times 10^{7} = 2\times
10^{6}$ बस्तियाँ; \qty{10}{\%}, यानी $2\times 10^{5}$, में अंतर्वेशन है।
\textbf{4.} सफ़ेद: \qty{10}{\%}। अंतर्वेशनों का आधा सही दिशा में है,
इसलिए $P(\text{कोई नहीं, } n \text{ सफ़ेदों में}) = 0.5^{n} \le 0.01$
से $n \ge 6.6$: सात उठाइए।
\textbf{5.} $32$ द्विगुणन, $2^{32} = 4.3\times 10^{9}$ कोशिकाएँ, $2.1\times
10^{11}$ \omterm{def:b3:genetic-engineering:tools}{प्लाज़्मिड}; हर एक $5500\times 650 = 3.6\times 10^{6}$ Da $=
5.9\times 10^{-18}$ g: \qty{1.3}{\micro g} \omterm{def:b3:genetic-engineering:tools}{प्लाज़्मिड}।
\textbf{6.} क्लोनन को केवल प्रतिकृति चाहिए, जो \omterm{def:b3:genetic-engineering:tools}{प्लाज़्मिड} का उद्गम देता
है; मानव प्रवर्तक को जीवाणु RNA पॉलीमरेज़ नहीं पढ़ती, इसलिए अभिव्यक्ति
के लिए (अंतर्वेशन-रहित) कूटलेखी अनुक्रम के आगे कोई जीवाणु प्रवर्तक और
राइबोसोम-बंधक स्थल रखना पड़ता है।
\textbf{7.} $N_{0} = N_{T}/2^{C_{T}}$: $10^{12}/2^{28} \approx 3700$
प्रतियाँ; $C_{T} = 35$ पर $10^{12}/2^{35} \approx 29$ प्रतियाँ।
\textbf{8.} $2^{n} = 10^{12}$: $n = 39.9$, चालीस चक्र।
\textbf{9.} $10^{12}/2^{38} \approx 3.6$ प्रतियाँ। लगभग $40$ चक्रों से
आगे कोई अकेला संदूषक अणु, या प्रारंभक-कृतक, देहली तक पहुँच जाते हैं, और
एक प्रति से कम वाला प्रतिदर्श निरर्थक है।
\textbf{10.} एक अणु चक्र $40$ पर पार करता है: एक धनात्मक। इसे लगाने और
उत्पाद सँभालने के लिए अलग कमरे तथा उपकरण, एकतरफ़ा कार्य-प्रवाह, हर बारी
में ऋणात्मक नियंत्रण, और ढुले हुए प्रवर्धित उत्पादों के एंज़ाइमी विनाश
से रोकिए।
\textbf{11.} $3700/0.4 \approx 9000$ RNA प्रतियाँ।
\textbf{12.} बताया गया $2^{3.3} \approx 10$ गुना; सही $1.9^{3.3} =
e^{3.3\ln 1.9} \approx 8.3$ गुना।
\textbf{13.} $40\times \qty{35}{\micro g} = \qty{1.4}{mg}$ प्रति दिन,
\qty{0.51}{g} प्रति वर्ष; $\times 10^{7}$: \qty{5100}{kg} प्रति वर्ष।
\textbf{14.} $5.1\times 10^{6}/5800 \approx 880$ mol; $5.3\times 10^{26}$
अणु।
\textbf{15.} $5.1\times 10^{6}$ g $/ (20$ g/L$) = 2.6\times 10^{5}$ L
$= \qty{260}{m^3}$ --- ताल का दसवाँ भाग।
\textbf{16.} $5100\times 8000 = 4.1\times 10^{7}$ kg अग्न्याशय, प्रति
अग्न्याशय \qty{0.1}{kg} की दर से: वर्ष में $4\times 10^{8}$ सूअर।
\textbf{17.} पुनर्योगज हॉर्मोन मानव हॉर्मोन के समान है, इसलिए वह कोई
प्रतिरक्षी नहीं उठाता और कोई एलर्जी नहीं करता; वह कोई पशु-रोगजनक
(विषाणु, प्रायॉन) नहीं ढोता; आपूर्ति वध पर निर्भर नहीं; और अनुक्रम
सुधारा जा सकता है।
\textbf{18.} जो अवशेष ग्राही को छूते हैं वे अपरिवर्तित हैं, इसलिए बंधन
और संकेतन इंसुलिन के ही हैं; बदले हुए अवशेष उस पृष्ठ पर हैं जहाँ
इंसुलिन अणु द्विलकों और षट्लकों में जुड़ते हैं, और उन्हें बदलने से यह
बदल जाता है कि अंतःक्षेपित भंडार कितनी तेज़ी से अवशोषणीय एकलकों में
टूटता है।
\textbf{19.} $6.4\times 10^{9}\times 4^{-20}/16 \approx 3.6\times
10^{-4}$: कोई ठीक-ठीक संयोगी मेल नहीं।
\textbf{20.} $1 + 60 + 1710 + 27\times 1140 = \num{32551}$ अनुक्रम;
$\num{32551}\times 3.6\times 10^{-4} \approx 12$ संयोगी लक्ष्येतर
स्थल।
\textbf{21.} $10^{7}\times 10^{-4} = 1000$ कोशिकाएँ। कोई स्तंभ कोशिका
रोगी के जीवन भर जीती और विभाजित होती है; जिसने कोई अर्बुद-दमनकारी खोया
है, या कोई स्थानांतरण पाया है, वह ऐसा क्लोन बसा सकती है जो कोई
श्वेतरक्तता बन जाए।
\textbf{22.} $p_{1} = 0.0019$, $p_{2} = 0.0036$, $p_{3} = 0.0069$;
प्रति पीढ़ी $1.9$ की गुणोत्तर वृद्धि से $\ln 500/\ln 1.9
\approx 10$ मिलता है; पुनरावृत्ति-संबंध दोहराने पर बारंबारता $0.5$
पीढ़ी $11$ पर पार होती है ($0.45 \to 0.68$)।
\textbf{23.} लगभग एक वर्ष। समयुग्मजियों में \qty{20}{\%} लागत
$c = 0.9$ से कम है, इसलिए चालक फिर भी स्थिरीकरण तक जाता है, और धीमे,
जबकि समष्टि की माध्य अनुकूलता गिरती है --- किसी दमन-चालक का यही
प्रयोजन है; $c$ से अधिक लागत उसे किसी मध्यवर्ती बारंबारता पर रोक देती।
\textbf{24.} एक पीढ़ी में $10^{6}\times 0.1 = 10^{5}$ प्रतिरोध
युग्मविकल्पी। वे कटने योग्य नहीं हैं और, यदि अनुकूल हों, तो चालक के
महँगे होते हुए भी निर्लागत हैं, इसलिए उन्हें वरीयता मिलती है और वे फैल
जाते हैं, और चालक समाप्त हो जाता है --- इसीलिए चालक ऐसे स्थल लक्ष्य
करते हैं जहाँ सिरा-संयोजन उत्पाद घातक हों, एक साथ कई मार्गदर्शकों के साथ।
\textbf{25.} पहले प्रतिदर्श में लगभग $3700$ प्रतियाँ; संसार भर के
इंसुलिन के लिए वर्ष में कोई \qty{260}{m^3} किण्वक; चालक के लिए लगभग
$11$ पीढ़ियाँ, यानी मच्छर का एक वर्ष।
\end{solution}
